\section{Conclusão}

A avaliação primária por validação cruzada agrupada mostrou AUC ROC de
$0{,}9709\pm0{,}0059$, aproximadamente $0{,}97\pm0{,}01$, para a ConvNeXt-Tiny na
triagem binária de células cervicais em citologia convencional. A variação da
sensibilidade entre 0,8253 e 0,9434 evidencia a importância de retreinar e
avaliar o modelo em múltiplas partições por imagem. Em uma partição retida, a
análise com TTA alcançou AUC ROC 0,9832 e permitiu explorar pontos
operacionais, incluindo sensibilidade 0,9701 com especificidade 0,8723.

A comparação local com ResNet50 indica vantagem da ConvNeXt-Tiny em
ranqueamento, mas não em todas as métricas dependentes de limiar. A análise
Bethesda aponta ASC-US como a principal dificuldade observada, com 81,7\% de
acerto entre 115 células na partição retida, resultado que deve ser confirmado
em folds e bases adicionais. A agenda futura inclui selecionar limiares por
fold e aplicar TTA na validação cruzada, repetir baselines sob o mesmo protocolo
agrupado, avaliar calibração com mais amostras, incorporar explicabilidade e
realizar validação externa. Também são necessários localização automática das
células e agregação das evidências no nível da lâmina. Até essas etapas, o
método permanece um protótipo de pesquisa para priorização, sem finalidade
diagnóstica.

\section*{Disponibilidade de Dados e Código}

A CRIC Cervix é de acesso público \cite{rezende2021cric}. Código e checkpoints
serão liberados em repositório público após a revisão.
